\documentclass[11pt]{article}
\usepackage[margin=1in]{geometry}
\usepackage{amsmath, amssymb, amsthm}
\usepackage{mathtools}
\usepackage{enumitem}
\usepackage{booktabs}
\usepackage{hyperref}
\usepackage{xcolor}
\hypersetup{colorlinks=true, linkcolor=blue!50!black, urlcolor=blue!50!black}

\newtheorem{theorem}{Theorem}
\newtheorem{proposition}[theorem]{Proposition}
\newtheorem{corollary}[theorem]{Corollary}
\newtheorem{lemma}[theorem]{Lemma}

\theoremstyle{definition}
\newtheorem{definition}[theorem]{Definition}
\newtheorem{remark}[theorem]{Remark}

\newcommand{\R}{\mathbb{R}}
\newcommand{\Rpos}{\R_{>0}}
\newcommand{\Rnn}{\R_{\geq 0}}
\newcommand{\bpmz}{\mathrm{BPM}_0}
\newcommand{\bpm}{\mathrm{BPM}}
\newcommand{\dB}{\,dB}
\newcommand{\dtau}{\,d\tau}
\newcommand{\dmu}{\,d\mu}
\newcommand{\Lone}{L^{1}}
\newcommand{\BV}{\mathrm{BV}}
\newcommand{\AC}{\mathrm{AC}}
\newcommand{\one}{\mathbf{1}}

\title{SV Engine: Mathematical Design}
\author{vsrg-analysis Player}
\date{}

\begin{document}
\maketitle

\begin{quote}
\textbf{TL;DR.} \emph{Time-space} is smooth: notes scroll at a velocity that's
a function of wall-clock time, and the cumulative position is a continuous
function. \emph{Beat-space} is not smooth, because charts can have
\emph{stops} (time passes, beats don't) and \emph{warps} (beats jump, time
doesn't) --- and a warp creates a literal jump in the cumulative-position
function that no time-space velocity can reproduce.

Came up by Yucky, written up by a \textbf{CLANKER}. If you want a even simpler explaination, 
think about it like the difference between an indestructible ball and a destructible
ball that is usually round but sometimes not (make a PR if you still don't understand this)

The rest of this document formalizes that observation: time-space is integration
against ordinary Lebesgue measure $d\tau$; beat-space is integration against a
\emph{Stieltjes measure} $dB$ that allows point masses (= warps). Once you
write it this way, ``which game's SV are we using'' collapses to ``which
measure are we integrating against,'' and switching engines at runtime is just
swapping $\mu$ on the chart document. One integrator, multiple measures, one
codebase.
\end{quote}

\begin{abstract}
This document gives the mathematical foundation for the unified scroll-velocity
(SV) framework used by the Player. The framework is built on a single
\emph{measure-theoretic primitive}: render distance is a Lebesgue--Stieltjes
integral of a velocity against a chart-derived measure on chart time. Engines
differ only in their choice of measure; the document carries the integrand and
the measure as orthogonal data.

Three structural facts drive the design. First, time-space and beat-space are
not different positioning theories but two choices of measure --- Lebesgue
$\dtau$ versus the Stieltjes measure $dB$ induced by the timing map.
Second, the failure of time-space to represent warps is a property of the
\emph{measure}: $\Lone(T, d\tau)$ functions cannot place mass at a point, so
no enrichment of the integrand alone can recover warps under Lebesgue measure.
Third, conversion between measures is exact wherever the timing map is
absolutely continuous, by Radon--Nikodym; warps are precisely the failure of
absolute continuity.

The engine architecture follows: the canonical document carries a measure $d\mu$
and an $\Lone(T, d\mu)$ integrand $v$. A single integrator computes
$\int v\dmu$. Engine swap reduces to swapping the measure on the document; the
integrator is unchanged. Per-column SV is $K$ integrands against one shared
measure; per-note SV is sparse-indexed integrands against the same measure.
\end{abstract}

\tableofcontents

% =====================================================================
\section{Function spaces and measures}
\label{sec:spaces}

We work on a closed chart-time interval $T = [t_0, t_{\mathrm{end}}] \subset \R$.
The framework is parameterized by a pair $(\mu, v)$ where $\mu$ is a positive
Radon measure on $T$ and $v$ is a $\mu$-integrable function. The render distance
is $\int v\dmu$. The two ``spaces'' the rhythm-game literature distinguishes
correspond to two choices of $\mu$.

\begin{definition}[Time-space]
\label{def:time-space}
The time-space framework is the pair $(d\tau, v)$ with $v \in \Lone(T, d\tau)$
and $d\tau$ Lebesgue measure on $T$.
\end{definition}

\begin{definition}[Beat-space]
\label{def:beat-space}
A \emph{timing map} is a non-decreasing right-continuous function
$B : T \to \Rnn$ of bounded variation, with $B(t_0) = 0$. The
\emph{Stieltjes measure} $dB$ is the unique Radon measure on $T$ with
\[
dB([\alpha, \beta]) \;=\; B(\beta) - B(\alpha^-)
\quad \text{for } \alpha \leq \beta \in T.
\]
The beat-space framework is the pair $(dB, s \circ B)$ with
$s : \Rnn \to \R$ measurable and $s \circ B \in \Lone(T, dB)$.
\end{definition}

\begin{remark}[Anatomy of $dB$]
\label{rem:dB-anatomy}
By the Lebesgue decomposition, every Radon measure on $T$ splits as
$dB = (dB)_{\AC} + (dB)_{\mathrm{sing}}$, where $(dB)_{\AC} \ll d\tau$ has a
Radon--Nikodym density $B'(\tau) \dtau$ and $(dB)_{\mathrm{sing}} \perp d\tau$
is supported on a Lebesgue-null set. For chart timing maps the singular part is
purely atomic: it is a finite sum of point masses, one at each warp:
\begin{equation}
\label{eq:dB-decomp}
dB \;=\; B'(\tau)\dtau \;+\; \sum_{w} \Delta b_w \cdot \delta_{\tau_w}.
\end{equation}
The atoms encode warps; the absolutely continuous part encodes BPM. Stops
contribute to $(dB)_{\AC}$ via $B' = 0$ on the stop interval --- they are
absolutely continuous, just degenerate.
\end{remark}

\begin{definition}[Position-dependent zoom]
The \emph{zoom} $z : T \to \Rpos$ is bounded measurable in chart time and is
applied multiplicatively at the playhead, outside the integral. Time-space
engines fix $z \equiv 1$; Etterna's \texttt{\#SPEEDS} produces a nontrivial
$z$.
\end{definition}

\begin{definition}[Base BPM]
$\bpmz \coloneqq \bpm(t_0)$ is the BPM of the first segment.
\end{definition}

\begin{definition}[Column index]
Let $C = \{0, \dots, K-1\}$ index the playfield columns. Where the integrand
varies per column we write $v(\tau, c)$ and $z(\tau, c)$. The measure $\mu$ is
column-independent: it is a property of the chart timing data only.
\end{definition}

% =====================================================================
\section{The render-distance functional}
\label{sec:functional}

\begin{definition}[Render distance]
\label{def:render-distance}
Given a measure $\mu$ on $T$, an integrand $v \in \Lone(T, \mu)$, and a zoom
$z$, the \emph{render distance} between chart times $a, b \in T$ as observed
from playhead $a$ is
\begin{equation}
\label{eq:d}
d_\mu(a, b) \;\coloneqq\; z(a) \cdot \int_{(a, b]} v(\tau)\dmu(\tau).
\end{equation}
We use half-open intervals $(a, b]$ throughout to make point-mass attribution
unambiguous: a warp at $\tau_w$ contributes to $d_\mu(a, b)$ iff $\tau_w \in
(a, b]$.
\end{definition}

\begin{definition}[Cumulative function]
\label{def:cumulative}
The \emph{cumulative function} is
\begin{equation}
\label{eq:C}
C_\mu(t) \;\coloneqq\; \int_{(t_0, t]} v(\tau)\dmu(\tau).
\end{equation}
\end{definition}

\begin{proposition}[Cumulative decomposition]
\label{prop:decomp}
For all $a, b \in T$,
\[
d_\mu(a, b) \;=\; z(a) \cdot \bigl(C_\mu(b) - C_\mu(a)\bigr).
\]
When $z \equiv 1$, $C_\mu$ alone determines render distance; cull-space queries
reduce to a single bisection on a precomputed array of $C_\mu$ samples.
\end{proposition}

\begin{remark}[Why $C$ is cached and $z$ is per-frame]
\label{rem:why-split}
$C_\mu$ depends only on the chart, so it is computed once at chart load and
queried at $O(\log n)$ per call. $z(a)$ depends on the playhead and is
evaluated per frame. Splitting them lets the cull-space window be a single
bisect over $C_\mu$, with $z$ applied as a uniform multiplier on the visible
delta. This is exactly the structure the existing \texttt{SVEngine} Protocol
exposes (\texttt{cumulative\_at}, \texttt{render\_multiplier\_at},
\texttt{distance}).
\end{remark}

\begin{remark}[Cull-space approximation under nontrivial $z$]
When $z$ varies, $z(a) \cdot (C_\mu(b) - C_\mu(a))$ approximates the true
integral by treating $z$ as constant over $(a, b]$. The error during a brief
$z$-transition is at most one segment's multiplier, which the Player pads for.
This matches Etterna's ArrowEffects.cpp XMOD branch.
\end{remark}

% =====================================================================
\section{Instantiations}
\label{sec:instances}

\subsection{Time-space (osu!mania, Quaver)}

$\mu = d\tau$. For osu!mania $z \equiv 1$:
\[
d_{d\tau}(a, b) \;=\; \int_{(a,b]} v(\tau)\dtau.
\]
The timing map $B$ does not appear. Stops and warps are not natively
expressible (\S\ref{sec:stops-warps}).

For Quaver $z(a) = \mathrm{ssf}(a)$ is the chart-authored
ScrollSpeedFactor, a piecewise-linear position-dependent zoom keyed at
chart load and lerped at the playhead. Quaver also authors $v$
per-group (\S\ref{sec:groups}); the per-group SSF means $z$ is per-group
too: $z(a, g) = \mathrm{ssf}_g(a)$.

\subsection{Beat-space (Etterna XMOD)}
\label{sec:beat-instance}

$\mu = dB$, integrand $v(\tau) = s(B(\tau))$, $z(a) = \mathrm{speed\_percent}(a)$.
By definition,
\begin{equation}
\label{eq:stieltjes}
d_{dB}(a, b)
\;=\; z(a) \cdot \int_{(a,b]} s(B(\tau))\dB(\tau).
\end{equation}
On any interval where $B$ is absolutely continuous, $dB = B'(\tau)\dtau$ and
\eqref{eq:stieltjes} reduces to a Lebesgue integral. Substituting $b = B(\tau)$
when $B$ is AC and strictly increasing,
\[
\int_{(a,b]} s(B(\tau)) B'(\tau)\dtau \;=\; \int_{B(a)}^{B(b)} s(\beta)\, d\beta
\;=\; D(B(b)) - D(B(a)),
\]
where $D(\beta) = \int_0^\beta s$ is the displayed-beat function. Converting to
render seconds via the base BPM:
\[
d_{dB}(a, b) \;=\; z(a) \cdot \frac{60}{\bpmz} \cdot \bigl(D(B(b)) - D(B(a))\bigr) \quad \text{(AC region)}.
\]
This is exactly Etterna's XMOD positioning equation. At a warp $\tau_w \in (a,b]$,
the singular part of $dB$ contributes an additional point-mass term
$s(b_w^+) \cdot \Delta b_w$ via \eqref{eq:dB-decomp}, which the AC formulation
above does not capture.

\subsection{Constant-BPM (IIDX)}

$\mu = dB$, $v \equiv 1$, $z \equiv h$ for player-chosen hi-speed $h$:
\[
d_{dB}(a, b) \;=\; h \cdot dB((a, b]) \;=\; h \cdot (B(b) - B(a)).
\]
Chart SV is dropped; BPM changes propagate through $B'$ and warps through the
atoms of $dB$, but the constant integrand neutralizes any chart-authored
multiplier. Floating hi-speed is a separate solve for $h$ at chart load and
lives outside the integrator.

\subsection{Identity}

$\mu = d\tau$, $v \equiv 1$, $z \equiv 1$: $d(a,b) = b - a$.

% =====================================================================
\section{Lossless conversion}
\label{sec:conv}

The conversion between time-space and beat-space is the Radon--Nikodym
relationship between the two measures, valid wherever it applies.

\begin{theorem}[Equivalence on AC intervals]
\label{thm:conversion}
Let $I \subseteq T$ be an interval on which $B$ is absolutely continuous (no
warps in $I$; stops are permitted). Let $s \in \Lone(B(I), d\beta)$ be a
beat-space integrand. Define
\begin{equation}
\label{eq:beat-to-time}
v^*(\tau) \;\coloneqq\; s(B(\tau)) \cdot \frac{\bpm(\tau)}{\bpmz}, \qquad \tau \in I.
\end{equation}
Then for all $a, b \in I$,
\[
\int_{(a,b]} v^*(\tau)\dtau
\;=\; \frac{60}{\bpmz} \int_{(a,b]} s(B(\tau))\dB(\tau).
\]
That is, the time-space integral of $v^*$ against $d\tau$ equals the
beat-space integral of $s \circ B$ against $dB$, up to the global factor
$60/\bpmz$.
\end{theorem}

\begin{proof}
On $I$, $B$ is AC, so by Radon--Nikodym $dB = B'(\tau)\dtau$ with
$B' \in \Lone(I, d\tau)$, and $B'(\tau) = \bpm(\tau)/60$ a.e. Then
\[
\int_{(a,b]} s(B(\tau))\dB(\tau)
\;=\; \int_{(a,b]} s(B(\tau))\, B'(\tau)\dtau
\;=\; \int_{(a,b]} s(B(\tau)) \cdot \frac{\bpm(\tau)}{60}\dtau.
\]
Multiplying by $60/\bpmz$ gives the right-hand side equal to
$\int_{(a,b]} s(B(\tau)) \cdot \bpm(\tau)/\bpmz \dtau = \int_{(a,b]} v^*(\tau)\dtau$.
\qedhere
\end{proof}

\begin{theorem}[Time $\to$ beat]
\label{thm:t2b}
Under the same hypotheses as Theorem~\ref{thm:conversion}, given a time-space
integrand $v \in \Lone(I, d\tau)$, define
\begin{equation}
\label{eq:time-to-beat}
s^*(\beta) \;\coloneqq\; v(B^{-1}(\beta)) \cdot \frac{\bpmz}{\bpm(B^{-1}(\beta))}
\qquad \text{for } dB\text{-a.e.\ } \beta \in B(I).
\end{equation}
Then $\int_{(a,b]} v(\tau)\dtau = (60/\bpmz) \int_{(a,b]} s^*(B(\tau))\dB(\tau)$.
\end{theorem}

\begin{proof}
Symmetric to Theorem~\ref{thm:conversion}. $B^{-1}$ is well-defined
$dB$-a.e.\ on $B(I)$: stops are atoms in the $\dtau$-pushforward but
$dB$-null in beat space (stops are intervals where $B$ is constant, so $dB$
assigns them measure zero), so the failure of pointwise invertibility on
stops occurs on a $dB$-null set and does not affect the integral.
\qedhere
\end{proof}

\begin{remark}[Where conversion fails]
\label{rem:conv-fails}
Theorem~\ref{thm:conversion} requires $B$ absolutely continuous on $I$.
Warps violate AC: at a warp $\tau_w$, $B$ has a jump
$B(\tau_w) - B(\tau_w^-) = \Delta b_w > 0$, so $dB$ has an atom
$\Delta b_w \cdot \delta_{\tau_w}$ that cannot be written as $B'(\tau)\dtau$
for any $B' \in \Lone$. The conversion theorem does not extend across warps;
\S\ref{sec:stops-warps} treats this as an irreducible obstruction.
\end{remark}

\begin{corollary}[Zoom is orthogonal]
\label{cor:z-ortho}
The zoom $z(a)$ is multiplied outside the integral and does not appear in
\eqref{eq:beat-to-time} or \eqref{eq:time-to-beat}. Conversion preserves $z$
verbatim; the engine independently chooses whether to honor it.
\end{corollary}

\begin{remark}[Distance preserved; cumulative not]
\label{rem:cum-not-invariant}
Theorems~\ref{thm:conversion} and \ref{thm:t2b} preserve $d_\mu(a, b)$ for all
$a, b$. The cumulative $C_\mu$ is \emph{not} preserved pointwise across a
change of measure: an antiderivative is measure-specific. Engine swap must
rebuild $C_\mu$ from scratch in the new measure; cached values are not
translatable.
\end{remark}

% =====================================================================
\section{Stops, delays, and warps}
\label{sec:stops-warps}

\subsection{Stops and delays}

A stop at chart time $\tau_s$ of duration $\delta$ has $B$ constant on
$[\tau_s, \tau_s + \delta]$. Then $dB$ assigns this interval measure zero, so
\[
\int_{(\tau_s, \tau_s + \delta]} s(B(\tau))\dB(\tau) \;=\; 0
\]
for any integrand --- notes do not move. In time-space, the analogous behavior
requires setting $v = 0$ on the stop and re-parameterizing the chart-time
domain by $+\delta$ for all subsequent events. This is a domain extension, not
a pointwise patch.

\subsection{Warps}

A warp at $\tau_w$ of beat extent $\Delta b > 0$ contributes the atom
$\Delta b \cdot \delta_{\tau_w}$ to $dB$. Then
\[
\int_{(a, b]} s(B(\tau))\dB(\tau)
\;=\; \int_{(a,b] \setminus \{\tau_w\}} s \cdot B'\dtau
\;+\; s(b_w^+) \cdot \Delta b \cdot \one_{\tau_w \in (a, b]}.
\]
The atom contributes a step jump in $C_{dB}$ at $\tau_w$ with size proportional
to the SCROLLS ratio at the warp's landing beat.

\begin{proposition}[Warp irreducibility]
\label{prop:warp-irreducible}
There is no $v \in \Lone(T, d\tau)$ such that
$\int_{(a, b]} v(\tau)\dtau$ has a jump discontinuity in $b$ at some
$\tau_w \in (a, b]$.
\end{proposition}

\begin{proof}
For $v \in \Lone(T, d\tau)$, the cumulative
$F(b) = \int_{(t_0, b]} v(\tau)\dtau$ is absolutely continuous in $b$, hence
continuous. No $\Lone$ density against Lebesgue measure produces jump
discontinuities in its antiderivative. \qedhere
\end{proof}

\begin{remark}[The obstruction is the measure, not the integrand]
\label{rem:obstruction-is-measure}
Proposition~\ref{prop:warp-irreducible} is stronger than ``you'd need a
Dirac.'' Even allowing $v$ to range over signed Radon measures rather than
$\Lone$ functions does not help \emph{within time-space}: if one extends the
integrand to a measure $v\dtau \to \nu$, then $\int_{(a,b]} \nu$ depends on
$\nu$, not on $\dtau$; one has effectively replaced the time-space pair
$(d\tau, v)$ with $(\nu, 1)$, which is a beat-space integral against a
different measure. Thus warps are not expressible by enriching the integrand
within the time-space framework: they require a different measure.
\end{remark}

\begin{corollary}[Canonical form]
\label{cor:canonical}
Any framework intended to support warps must integrate against a measure with
nontrivial singular part. The canonical \texttt{SVData} stores $dB$
with explicit warp atoms, making beat-space the structurally complete
representation. Time-space-authored charts (osu!, Quaver) embed by setting the
singular part to zero.
\end{corollary}

% =====================================================================
\section{BPM--SV coupling toggle}
\label{sec:coupling}

Some authoring formats (Quaver) expose a flag controlling whether timing-point
BPM scales SVs (\texttt{BPMDoesNotAffectScrollVelocity}). In the present
framework this is not a separate space: it is a choice of which integrand
the chart authored.

\begin{definition}[Coupled and decoupled time-space]
The \emph{BPM-coupled} time-space integral is
$\int v(\tau) \cdot \bpm(\tau)/\bpmz \dtau$;
the \emph{BPM-decoupled} integral is the plain $\int v(\tau)\dtau$.
\end{definition}

\begin{proposition}[Coupling is a choice of authoring frame]
\label{prop:coupling}
A BPM-coupled time-space integrand $v_{\mathrm{coup}}$ produces the same
render distance as a beat-space integrand $s$ on AC intervals iff
$v_{\mathrm{coup}} = s \circ B$ (after multiplying by $60/\bpmz$).
\end{proposition}

\begin{proof}
Direct from Theorem~\ref{thm:conversion} with $v^* = (s \circ B) \cdot \bpm/\bpmz$.
\qedhere
\end{proof}

The flag therefore declares which authoring frame the file uses; the canonical
document stores the equivalent decoupled form, and the engine applies coupling
at evaluation if requested by the chart's flag.

% =====================================================================
\section{Per-column and per-note SV}
\label{sec:per-column}

\begin{definition}[Selector]
A \emph{selector} is an element of
$\mathrm{All} \,\sqcup\, 2^C \,\sqcup\, 2^{\mathcal{N}}$
($\mathrm{All}$, a column subset, or a note-id subset).
\end{definition}

Each velocity record carries a selector. The per-column integrand is
\[
v(\tau, c) \;=\; m_k \quad \text{where } k = \max\{ i : \tau_i \leq \tau,\ c \in \mathrm{sel}(i)\}.
\]

\begin{theorem}[Per-column independence]
\label{thm:per-column}
Fix a measure $\mu$ on $T$ and a per-column integrand $v(\tau, c)$ with
$v(\cdot, c) \in \Lone(T, \mu)$ for each $c$. Define
$d_c(a, b) = z(a, c) \cdot \int_{(a,b]} v(\tau, c)\dmu(\tau)$ and
$C_c(t) = \int_{(t_0, t]} v(\tau, c)\dmu(\tau)$. Then
\S\ref{sec:functional}--\S\ref{sec:conv} apply column-wise: the conversion of
Theorem~\ref{thm:conversion} preserves $d_c$ for each $c$ independently,
and stop/warp behavior under $\mu = dB$ is per-column identical.
\end{theorem}

\begin{proof}
$\mu$ is column-independent and the analyses of
\S\ref{sec:functional}--\S\ref{sec:conv} depend on $v$ only through its
$\Lone(T, \mu)$ membership. Apply column-wise. \qedhere
\end{proof}

\begin{remark}[Storage and per-note]
The cumulative cache becomes $K$ parallel arrays $C_0, \dots, C_{K-1}$.
Records with selector $\mathrm{All}$ share segments structurally; per-column
records branch the affected columns. Per-note SV (selector
$N \subseteq \mathcal{N}$) is type-compatible: the integrand becomes
$v(\tau, n)$ indexed by note. Sparse storage is required for efficient
$|\mathcal{N}|$-many cumulative arrays; the math (Theorem~\ref{thm:per-column})
extends without change.
\end{remark}

\subsection{Group selectors (Quaver TimingGroups)}
\label{sec:groups}

A common selector pattern is partition by \emph{group}: the note set
$\mathcal{N}$ is partitioned into disjoint groups $\{G_g\}_g$, each
group $g$ owns a stream $v_g \in \Lone(T, \mu)$, and a note $n \in G_g$
inherits $v(\tau, n) = v_g(\tau)$. Quaver authors this directly via the
\texttt{TimingGroup} field on each hit object; the chart's top-level
\texttt{SliderVelocities} populate a distinguished default group
$g = \mathrm{def}$.

The group partition collapses storage: instead of one cumulative cache
per note, we keep one $C_g$ per group and dispatch by group at lookup.
Each $C_g$ is built independently against the same shared $\mu$, so
Theorem~\ref{thm:per-column} applies to each group in isolation
(per-column independence is a special case of per-group independence
under a column-set selector).

\begin{remark}[Playhead is per-group]
\label{rem:playhead-per-group}
Render distance from a playhead $a$ to a note $n \in G_g$ is
$d_g(a, t_n) = z(a, g) \cdot (C_g(t_n) - C_g(a))$. The playhead's
``current cumulative'' is therefore \emph{per group}: $C_g(a)$ for the
note's own group, not $C_{\mathrm{def}}(a)$. Mixing the two -- e.g.,
subtracting the default-group playhead from a non-default-group note
position -- produces the wrong distance whenever $v_g$ diverges from
$v_{\mathrm{def}}$. The implementation must evaluate the playhead in
each group's stream.
\end{remark}

% =====================================================================
\section{Extended notes: long notes under reversal}
\label{sec:long-notes}

A long note (LN) covers a chart-time interval $[t_h, t_e]$. Its body
must be drawn as a continuous bar in render-space; under signed
integrands the cumulative function need not be monotonic, so head and
tail cumulative values do not bound the body.

\begin{definition}[LN body extent]
\label{def:body-extent}
For $n \in G_g$ a long note with head $t_h$ and end $t_e$, the
\emph{body extent} of $n$ is the smallest interval enclosing $C_g$
restricted to $[t_h, t_e]$:
\begin{equation}
\label{eq:body-extent}
\mathcal{B}_g(t_h, t_e)
\;=\; \Bigl[\, \min_{t \in [t_h, t_e]} C_g(t),\ \max_{t \in [t_h, t_e]} C_g(t)\,\Bigr].
\end{equation}
The body's screen-space rectangle is the image of $\mathcal{B}_g(t_h, t_e)$
under $C_g \mapsto z(a, g) \cdot (C_g - C_g(a))$ at playhead $a$.
\end{definition}

\begin{proposition}[Extrema at sign changes]
\label{prop:body-extrema}
Let $v_g$ be piecewise-constant with breakpoints
$t_h = \tau_0 < \tau_1 < \cdots < \tau_K = t_e$ on $[t_h, t_e]$, and
let $\mu = d\tau$ on this interval (no atoms, time-space). Then
\begin{equation}
\label{eq:hull-vertices}
\mathcal{B}_g(t_h, t_e)
\;=\; \bigl[\min S,\ \max S\bigr],
\end{equation}
where
\[
S \;=\; \{C_g(t_h),\, C_g(t_e)\}
\;\cup\; \{ C_g(\tau_k) : 1 \leq k < K,\ \mathrm{sign}(v_g(\tau_k^+)) \neq \mathrm{sign}(v_g(\tau_k^-)) \}.
\]
That is, the body extent is determined by the endpoints together with
the cumulative values at sign-change breakpoints.
\end{proposition}

\begin{proof}
On each subinterval $(\tau_k, \tau_{k+1})$ the integrand $v_g$ is
constant, so $C_g$ is linear (slope $v_g(\tau_k^+) \cdot \rho_k$) and
attains its extrema on this subinterval only at the endpoints.
Therefore $\min$ and $\max$ of $C_g$ on $[t_h, t_e]$ occur at some
$\tau_k$. Suppose $\mathrm{sign}(v_g(\tau_k^+)) = \mathrm{sign}(v_g(\tau_k^-))$,
i.e., the slope sign does not change at $\tau_k$. Then $C_g$ is
monotone on $(\tau_{k-1}, \tau_{k+1})$, so $\tau_k$ is not a local
extremum of $C_g \restriction [\tau_{k-1}, \tau_{k+1}]$; in particular
it is dominated by one of $C_g(\tau_{k-1})$, $C_g(\tau_{k+1})$ in the
direction of the global extremum. Iterating eliminates every interior
non-sign-change breakpoint, leaving only $S$ as candidate vertices.
\qedhere
\end{proof}

\begin{remark}[Algorithm and storage]
Proposition~\ref{prop:body-extrema} reduces body-extent computation to
one $\Lone$-style scan: bisect $v_g$'s breakpoint array to locate
$[t_h, t_e]$, walk the slice, and emit $C_g(\tau_k)$ at each sign flip.
The cost per LN is $O(\#\text{breakpoints in } [t_h, t_e])$ at chart
load, with body extents stored as two scalars per LN
$(\mathcal{B}^{\min}_n, \mathcal{B}^{\max}_n)$. No per-frame work is
needed beyond the $C_g(a) \mapsto y$ projection that draws every other
note.
\end{remark}

\begin{remark}[Atom-supporting measures]
\label{rem:body-atoms}
Under $\mu = dB$ with warp atoms in $[t_h, t_e]$, $C_g$ has jump
discontinuities at the warp times $\tau_w$. The set $S$ in
Proposition~\ref{prop:body-extrema} extends to include
$C_g(\tau_w^-)$ and $C_g(\tau_w^+)$ for each $\tau_w \in (t_h, t_e)$;
the proof carries over because $C_g$ is monotone on each open
subinterval between consecutive breakpoints (and a warp atom is, by
the Lebesgue decomposition of \eqref{eq:dB-decomp}, a single jump
between two AC pieces). No game currently authors warps inside an LN
body, so the implementation handles only the $\mu = d\tau$ case.
\end{remark}

\subsection{Tail orientation}
\label{sec:tail-flip}

The LN tail sprite has a direction (head-side vs tail-side cap). Under
signed $v_g$, the tail can render below the head (forward scroll),
above it (reverse scroll at the tail), or anywhere in between. The
correct sprite orientation is determined by the sign of the integrand
at the tail.

\begin{definition}[Tail flip]
\label{def:tail-flip}
For $n \in G_g$ with end time $t_e$, the tail orientation flips when
$v_g(t_e^-) < 0$, where $v_g(t_e^-)$ is the most recent non-zero value
of $v_g$ at or before $t_e$. (Zero-multiplier segments don't change
the rendering direction; we walk backward through them.)
\end{definition}

\begin{remark}[Mirrors Quaver]
Quaver implements this as
\texttt{ScrollGroupController.IsSVNegative(EndTime)} and applies the
flip at draw time via a vertical sprite mirror. In our framework the
flip is a property of $v_g$ alone -- $z$ does not enter (it can scale
the body but not invert direction in any game we model).
\end{remark}

% =====================================================================
\section{Numerical exactness}
\label{sec:numerics}

\begin{proposition}[Exactness for piecewise-constant integrands]
\label{prop:exact}
If $v$ is piecewise-constant on a finite partition of $T$ and $\mu$ is a finite
sum of an AC part with piecewise-constant density plus finitely many atoms,
then $C_\mu$ is piecewise-linear-with-jumps on the common refinement of
the partitions. Evaluation of $C_\mu(t)$ at any $t$ reduces to one bisection
plus a constant-time arithmetic expression
\[
C_\mu(t) \;=\; C_\mu(\tau_k) + (t - \tau_k) \cdot v_k \cdot \rho_k
\;+\; \sum_{w : \tau_w \in (\tau_k, t]} v(\tau_w^+) \cdot a_w,
\]
where $\rho_k$ is the AC density of $\mu$ on segment $k$ and $a_w$ is the
atomic mass at warp $\tau_w$. No quadrature is required.
\end{proposition}

\begin{remark}[Test invariant]
For any $\mathrm{SVData}$ and any sample time $t$, the integrator's
$C_\mu(t)$ must agree with the brute-force reference
\[
C_{\mathrm{ref}}(t) \;=\; \sum_i v_i \cdot \mu(I_i \cap (t_0, t])
\]
to numerical precision (modulo summation order). This is a property test
across the full $\mathrm{SVData}$ space.
\end{remark}

\begin{remark}[Eased segments]
Quaver SSF and fluXis multiplier events relax piecewise-constant $v$ to
piecewise-smooth (linear, sinusoidal, etc.). The framework is unchanged: the
integrator must compute per-segment closed-form antiderivatives instead of
the constant case in Proposition~\ref{prop:exact}. Numerical exactness becomes
exactness up to the analytic accuracy of the easing-curve antiderivative.
\end{remark}

% =====================================================================
\section{Architectural mapping}
\label{sec:arch}

The structural payoff of the measure-theoretic framing is that the engine
collapses to a single integrator parameterized by the document's measure.
Engine identity reduces to a choice of $\mu$.

\subsection{The integrator}

A single procedure handles all engines:
\[
\texttt{cumulative\_at}(t, \mathrm{doc}) \;=\;
\int_{(t_0, t]} v(\tau)\, d\mu_{\mathrm{doc}}(\tau).
\]
The procedure decomposes the integral by segments of the common refinement
(Proposition~\ref{prop:exact}) and accumulates AC contributions and atoms
separately. \texttt{distance}, \texttt{render\_multiplier\_at}, and
\texttt{inverse\_cumulative\_at} are derived from this primitive.

\subsection{SVData}

The canonical document carries:
\begin{itemize}[leftmargin=2em]
    \item $\mu$: the integration measure, stored as
          $(\rho(\tau)\dtau, \{(\tau_w, a_w)\})$ --- a piecewise-constant AC
          density plus a list of atoms. For time-space-authored charts,
          $\rho \equiv 1$ and the atom list is empty. For beat-space-authored
          charts, $\rho = B'$ and atoms are warps.
    \item Velocity records: piecewise-constant $v$ with selectors per
          \S\ref{sec:per-column}. Each record carries $(\tau, b, m, \mathrm{sel}, \mathrm{group})$;
          the engine uses $\tau$ for time-space evaluation, $b$ for beat-space.
    \item Speed curves $z$: position-dependent zoom, optional, with
          per-segment easing.
    \item Time offsets: per-(time, group) animated visual offsets, optional.
    \item Coupling flag: BPM-coupled vs.\ BPM-decoupled
          (\S\ref{sec:coupling}).
\end{itemize}

By Corollary~\ref{cor:canonical}, the canonical form supports atoms in $\mu$;
time-space-authored charts embed losslessly by leaving the atom list empty.

\subsection{Engines as measure choices}

Each engine corresponds to a fixed $\mu$:

\begin{center}
\begin{tabular}{lccccc}
\toprule
Engine & $\mu$ & atoms & $z$ & warps & per-col \\
\midrule
TimeSpace      & $d\tau$           & dropped  & opt   & dropped   & opt \\
BeatSpace      & $dB$              & honored  & yes   & native    & opt \\
ConstantBPM    & $dB$, $v\equiv 1$ & honored  & const & dropped   & no \\
Identity       & $d\tau$, $v\equiv 1$ & ---   & no    & ---       & no \\
\bottomrule
\end{tabular}
\end{center}

The TimeSpace and BeatSpace engines run \emph{the same integrator} on the same
$\mathrm{SVData}$; they differ only in the $\mu$ they request. ConstantBPM
additionally pins $v \equiv 1$ and $z$ to the player's hi-speed. Identity
pins both to constants.

\subsection{Engine swap}
\label{sec:swap}

Per Remark~\ref{rem:cum-not-invariant}, swap reduces to:
\begin{enumerate}[leftmargin=2em]
    \item Reconfigure the integrator's measure to the new engine's $\mu$.
    \item Discard $C_\mu$ and rebuild from the document.
    \item Re-anchor the clock smoother by mapping the current chart-time $t$
          forward through the new $C_\mu$, not by translating the old value.
    \item Refresh \texttt{sv\_enabled} from the new engine's flag; if a
          CMOD-style override is suspended, refresh that state too.
\end{enumerate}
The renderer's per-frame contract is unchanged.

% =====================================================================
\section{Per-frame predictor (cull-space)}
\label{sec:predictor}

The renderer queries $C(t)$ every frame to position notes in cull-space.
Two correctness criteria must be met:

\begin{description}
    \item[Exactness on audio callback.] When the audio engine delivers a
        fresh chart-time $t_{\mathrm{raw}}$, the rendered cumulative must
        equal $C(t_{\mathrm{raw}})$.
    \item[Smoothness between callbacks.] Between callbacks, the rendered
        cumulative must advance linearly in wall-time at the local rate
        $r(t) = v(t)\,w(t)$, without observing per-frame noise from the
        audio engine's source clock.
\end{description}

The naive renderer evaluates $C(t_{\mathrm{raw}})$ per frame. This
satisfies exactness but fails smoothness: $t_{\mathrm{raw}}$ inherits
sub-callback stream-clock jitter on every per-frame read, producing
output position noise of order $r \cdot \sigma_{\mathrm{clock}}$, visible
at high local rates.

\subsection{The anchor-and-extrapolate algorithm}

Maintain a four-tuple of state:
\[
\bigl(t_a,\; C_a,\; r_a,\; W_a\bigr)
\]
where $t_a$ is the chart-time at the most recent anchor,
$C_a = C(t_a)$ exactly, $r_a = r(t_a^+)$ from the engine's
\texttt{cumulative\_velocity\_at}, and $W_a$ is wall-time
($\mathtt{monotonic()}$) at the anchor instant.

\begin{definition}[Per-frame query]
\label{def:query}
At wall-time $W$ with playback rate $\rho$ (chart-seconds per
wall-second), the predicted cumulative is
\begin{equation}
\label{eq:predict}
\widetilde{C}(W) \;=\; C_a + r_a \cdot (W - W_a) \cdot \rho.
\end{equation}
\end{definition}

\begin{theorem}[Equivalence on a constant-rate segment]
\label{thm:predictor-equiv}
On a segment where $r$ is constant on $[t_a, t_a + (W - W_a)\rho]$ and no
warp atom intervenes, $\widetilde{C}(W) = C(t_a + (W - W_a)\rho)$ exactly.
\end{theorem}

\begin{proof}
On a constant-rate interval, $C(t) - C(t_a) = \int_{t_a}^{t} r\,d\tau =
r \cdot (t - t_a)$. Substituting $t = t_a + (W - W_a)\rho$ gives
$C(t) = C_a + r \cdot (W - W_a)\rho$, which equals \eqref{eq:predict}.
\qedhere
\end{proof}

\subsection{Re-anchoring on breakpoints}

When the predicted chart-time $t_a + (W - W_a)\rho$ would cross the next
breakpoint $b$, the algorithm advances the anchor to that breakpoint:
\[
\bigl(t_a',\;C_a',\;r_a',\;W_a'\bigr) \;=\;
\bigl(b,\; C(b),\; r(b^+),\; W_a + (b - t_a)/\rho\bigr).
\]
$C(b)$ comes from the integrator (exact, including any warp atom at $b$);
$r(b^+)$ is the post-boundary local rate. Subsequent per-frame queries
extrapolate from this new anchor. Multiple breakpoints in one frame are
handled by iterating until the predicted chart-time falls within the
current cell.

\begin{theorem}[Equivalence across breakpoints]
\label{thm:predictor-bp}
For any $W$ and any sequence of breakpoints $b_1 < b_2 < \dots < b_n$ in
$(t_a, t_a + (W - W_a)\rho]$, after $n$ re-anchor iterations,
$\widetilde{C}(W) = C(t_a + (W - W_a)\rho)$.
\end{theorem}

\begin{proof}
By induction on $n$. The base case $n = 0$ is Theorem~\ref{thm:predictor-equiv}.
Inductive step: after re-anchoring at $b_n$, the predictor's state is
$(b_n, C(b_n), r(b_n^+), W_a + (b_n - t_a)/\rho)$. The remaining
extrapolation crosses no further breakpoints in the current frame
(otherwise the loop continues), so by Theorem~\ref{thm:predictor-equiv}
the final value is $C(b_n) + r(b_n^+) \cdot (\text{remaining chart-time}) =
C(t_{\mathrm{end}})$, where the second equality holds because the
integrator splits its grid at $b_n$ and integrates exactly across each
sub-segment. \qedhere
\end{proof}

\subsection{Re-anchoring on audio callbacks and discontinuities}

Audio callbacks update $t_{\mathrm{raw}}$ ahead of the predictor's
extrapolation. The predictor detects this by comparing
$t_{\mathrm{raw}} - t_a$ against the predicted chart-time advance
$(W - W_a)\rho$:
\[
\bigl|\,(t_{\mathrm{raw}} - t_a) \;-\; (W - W_a)\rho\,\bigr| > \theta
\;\Longrightarrow\;\text{re-anchor at } t_{\mathrm{raw}}.
\]
The threshold $\theta$ is chosen larger than typical sub-callback clock
jitter ($\theta = 5\,\mathrm{ms}$ in the implementation) so per-frame
stream-clock noise does not trigger spurious re-anchors. Backward jumps
in $t_{\mathrm{raw}}$ (seek, scrub release) trigger immediate re-anchor;
explicit \texttt{reset(t)} calls (engine swap, rate change) re-anchor
unconditionally.

\subsection{Eased segments}

For an eased cell where $v(\tau)$ is linear from $v_s$ at $\tau = b_k$ to
$v_e$ at $\tau = b_{k+1}$, the integrator computes the per-cell
contribution exactly via the trapezoid rule:
\[
\int_{b_k}^{b_{k+1}} v\,w\,d\tau \;=\; \tfrac{1}{2}(v_s + v_e) \cdot w
\cdot (b_{k+1} - b_k).
\]
Within the cell, the integrator's $C(t)$ is quadratic in
$(t - b_k)$. The predictor's linear extrapolation at $r_a = v_s \cdot w$
(start-of-cell rate) is therefore inexact within the cell, with
worst-case error
\[
|\,\widetilde{C}(t) - C(t)\,| \;\leq\; \tfrac{1}{2}|v_e - v_s| \cdot w
\cdot (b_{k+1} - b_k),
\]
attained at the cell midpoint. The error vanishes at the cell endpoints,
where the predictor re-anchors. This is acceptable in practice because
authored eased cells are short ($\sim$ms scale for SSF / fluXis events),
and re-anchor at the cell endpoint zeroes the error before it propagates.

\subsection{Performance and jitter trade}

Measured with $50$ SV breakpoints across $600\,\mathrm{s}$ of chart, in
the Python implementation:
\begin{itemize}[leftmargin=2em]
    \item Per-frame $C(t_{\mathrm{raw}})$: $\sim 1.0\,\mu s$.
    \item Predictor: $\sim 2.9\,\mu s$.
\end{itemize}
At $60\,\mathrm{Hz}$, the predictor adds $\sim 3\,\mu s/$frame --
$0.018\%$ of the frame budget.

Jitter reduction, measured by frame-to-frame
$\Delta C$ standard deviation under Gaussian noise on $t_{\mathrm{raw}}$
with constant rate:
\begin{center}
\begin{tabular}{rccc}
\toprule
$\sigma_{\mathrm{in}}$ (ms RMS) & 0.10 & 0.50 & 1.00\\
\midrule
$\Delta C$ stddev (no predictor) & 0.000435 & 0.001448 & 0.002841\\
$\Delta C$ stddev (predictor) & 0.000331 & 0.000320 & 0.000573\\
ratio & $1.3\times$ & $4.5\times$ & $5.0\times$\\
\bottomrule
\end{tabular}
\end{center}
In the realistic regime ($\sigma_{\mathrm{in}} \approx 0.5$--$1\,\mathrm{ms}$
for PortAudio's \texttt{stream.time}), the predictor reduces visible
jitter by $\sim 5\times$. At $60\,\mathrm{Hz}$ with
$\mathrm{scroll\_speed} = 500\,\mathrm{px/s}$, this turns $\sim 1$ pixel
of frame-to-frame jitter into $\sim 0.16$ pixels -- below visual
perception threshold for most viewers.

\subsection{Native-port boundary}

The predictor and integrator are pure-functional given their inputs (no
Player-state coupling, no thread state). The fast path of
$\widetilde{C}$ in native code is one
$\mathtt{monotonic()}$ syscall, $\sim 4$ floating-point operations, and
two branches -- estimated $\sim 10$--$20\,\mathrm{ns}$ per frame in
$\mathrm{C}$ / Rust / Zig. The Python implementation's $\sim 3\,\mu s$
overhead is interpreter cost; lifting the math to a native module
behind an FFI shim would make the perf trade fully favorable.

The state struct, breakpoint cursor, and re-anchor protocol map to a
native module without re-deriving anything; the boundary is marked in
the source files.

% =====================================================================
\section{Open extensions}

\begin{itemize}[leftmargin=2em]
    \item \textbf{Per-note SV}: implemented for Quaver via group
          selectors (\S\ref{sec:groups}). Per-note selectors that don't
          partition cleanly (one-off SV per note) are still
          unimplemented; the math (Theorem~\ref{thm:per-column})
          extends without change but storage layout would differ.
    \item \textbf{LN body extent under reversal}: implemented per
          Proposition~\ref{prop:body-extrema} for $\mu = d\tau$. The
          $\mu = dB$ extension (warps inside an LN body,
          Remark~\ref{rem:body-atoms}) is unimplemented; no game
          authors this configuration.
    \item \textbf{Modscripting}: adds a per-frame mutation of $v$, $z$, and
          $\mathrm{sel}$. Out of scope; the framework's static-positioning
          subset is sufficient for non-modfile charts.
    \item \textbf{Eased SV segments}: Proposition~\ref{prop:exact} relaxes to
          per-segment closed-form antiderivatives; Remark on numerical
          exactness applies. Quaver SSF is implemented as an eased
          \emph{zoom} curve $z$ (linear lerp between keyframes); no
          eased \emph{integrand} segments yet.
    \item \textbf{Singular non-atomic measures}: the framework permits
          $\mu_{\mathrm{sing}}$ that is non-atomic (e.g.\ Cantor-like).
          No game uses these; the framework would handle them in principle
          but no engine implements such measures.
\end{itemize}

\end{document}
